\chapter*{Apéndice}
\addcontentsline{toc}{chapter}{Apéndice}

\setcounter{section}{0}
\renewcommand{\thesection}{A.\arabic{section}}

\section{Funciones holomorfas de varias variables}

No es objetivo indagar demasiado en el estudio de las funciones complejas en varias variables. Sin embargo, es necesario presentar el concepto de holomorfia para definir una variedad compleja. La mayoria de resultados se pueden encontrar en \parencite[Cap 1, Sec 1.1]{huybrechts2005complex}\\
Fijemos $U\subset\cc^n$ un subconjunto abierto. Para cualquier $w\in U$, definimos el \textbf{polidisco} como el conjunto definido por
\begin{equation*}
	B_\epsilon(w):=\{z: |z_i-w_i|<\epsilon_i\},\enspace \text{donde}\enspace \epsilon=(\epsilon_1,\epsilon_2,\cdots,\epsilon_n)
\end{equation*}

\begin{definition}[Función holomorfa]\label{taylor}
	Una función $f:U\subset \cc^n\to\cc$ es llamada holomorfa en un punto $w\in U$, si existe un polidisco $B_\epsilon(w)\subset U$ tal que la restricción $f\left.\right|_{B_\epsilon(w)}$ está dada por una serie de potencias convergente \begin{equation*}
		\sum_{(i_1,i_2,\cdots,i_n)\in\nn^n}^{\infty}a_{i_1,i_2,\cdots,i_n}(z_1-w_1)^{i_1}(z_2-w_2)^{i_2}\dots(z_n-w_n)^{i_n}
	\end{equation*}
	y es única.
\end{definition}

\begin{remark}[Definiciones equivalentes de holomorfia] Sea $f:U\to\cc$ una función. Entonces $f$ es holomorfa si 
	\begin{enumerate}
		\item Satisface las ecuaciones de Cauchy-Riemann en todas las coordenadas $z_j=x_j+iy_j$, es decir,
		\begin{equation*}
			\frac{\partial f}{\partial \overline{z_j}}=0,\qquad j=1,2,\cdots n
		\end{equation*}
		donde $\frac{\partial}{\partial \overline{z_j}}=\frac{1}{2}(\frac{\partial}{x}+i\frac{\partial }{\partial y})$.
		
		\item Es continuamente diferenciable y para todo $z\in U$ se cumple la siguiente fórmula
		\begin{equation*}
			f(z)=\frac{1}{(2\pi i)^n}\int_{\partial B_\epsilon(w)}\frac{f(w)}{(z_1-w_1)\dots(z_n-w_n)}dw_1\cdots dw_n
		\end{equation*}
	\end{enumerate}
	
\end{remark}

Veamos las generalizaciones de los resultados básico en una variable.

\begin{theorem}[Principio del máximo]\label{Principiodelmodulomaximo}
	Sea $U\subset \cc^n$ un conjunto abierto y conexo. Si $f:U\to\cc$ es una función holomorfa no constante, entonces $|f|$ no tiene un máximo local en $U$.
\end{theorem}

\begin{theorem}[de la identidad]\label{identidad}
	Sea $U\subset \cc^n$ un subconjunto abierto y conexo. Si $f,g:U\to \cc$ son funciones holomorfas tales que $f(z)=g(z)$ para todo $z$ es un abierto no vacio de $U$, entonces $f=g$.
\end{theorem}

\begin{theorem}[de Lioville]
	Toda función $f:\cc^n\to\cc$ holomorfa acotada es constante.  
\end{theorem}

\begin{theorem}[Hartogs]\label{hartogs}
	Sea $\Omega\subset\cc^n$ con $n\geq 2$ un dominio y $K\subset\Omega$ un subconjunto compacto tal que $\Omega\setminus K$ es también conexo. Entonces, toda función holomorfa $f:\Omega\setminus K\to\cc$ se extiende holomorfamente a todo $\Omega.$
\end{theorem}


\begin{definition}\label{mapholocn}
	Una función $f:U\to\cc^m$ es llamada \textbf{aplicación holomorfa} si todas las funciones coordenadas $f_1,\dots,f_m:U\to\cc$ son funciones holomorfas.
\end{definition}


\begin{definition}
	Sean $U,V\subset \cc^n$ dos abiertos no vacios. Un aplicación holomorfa $f:U\to V$ es \textbf{biholomorfa} si es biyectiva y su inversa $f^{-1}$ también es una aplicación holomorfa. Además decimos que $U$ y $V$ son biholomorfos. 
\end{definition}

El caso del teorema del mapeo de Riemann es un caso que no se puede extender a más dimensiones y es exclusivo para $n=1$.
\begin{theorem}
	Sea $B_1^n(0)=\{z\in\cc^n: \|z\|<1\}$ la bola euclidiana unitaria y $P_1^n(0)=\{z\in\cc^n:|z_i|<1,\text{ para cada }i=1,2,\cdots, n\}$ el polidisco unitario en $\cc^n$. Entonces $B_1^n(0)$ y $P_1^n(0)$ son biholomorfos si, y solo si, $n=1$. La demostración se encuentra en \parencite[Teo 3.9, Pg. 78]{scheidemann}
\end{theorem}


\begin{definition}\label{puntoregular}
	Sean $U\subset \cc^m$ un conjunto abierto y $f:U\to\cc^n$ una aplicación holomorfa. El \textbf{jacobiano} (complejo) de $f$ en un punto $z\in U$ es la matriz 
	\begin{equation*}
		J(f)(z):=\left(\frac{\partial f_i}{\partial z_j}(z) \right)_{i,j}
	\end{equation*}
	Un punto $z\in U$ es llamado \textbf{regular} si el rango de la matriz jacobiana $J(f)(z)$ es $n$. Si todo punto $z\in f^{-1}(w)$ es regular, decimos que $w$ es un \textbf{valor regular}.
\end{definition}



\begin{theorem}[Teorema de la función inversa]
	Sea $f:U\to V$ una aplicación holomorfa entre dos subconjuntos abiertos $U,V\subset \cc^n$ y $z\in U$. Entonces los siguiente enunciados son equivalentes:
	\begin{enumerate}
		\item $\det J(f)(z)\neq 0$
		\item Existen vecindades $z\in U'\subset U$ y $f(z)\in V'\subset V$ tal que 
		\begin{equation*}
			f\left.\right|_{U'}:U'\to V'
		\end{equation*}
		es un biholomorfismo.
	\end{enumerate}
\end{theorem}

\begin{theorem}[Teorema de la función implícita]\label{funcionimplicita}
	Sea $U\subset\cc^{n+m}$ un abierto y $f:U\to \cc^n$ una aplicación holomorfa en $U$. Si $(w_0,z_0)\in U$ es un punto tal que
	$$ \det\left(\frac{\partial f_i}{\partial z_j}(w_0,z_0)\right)_{1\leq i,j \leq n}\neq 0$$
	Entonces existe una vecindad $W\subset \cc^m$ de $w_0$, una vecindad de $Z\subset\cc^n$ de $z_0$ y una aplicación holomorfa $g:W\to Z$ tal que $f(w,g(w))=0$ para todo $w\in W$ y
	\begin{equation*}
		\{(w,z)\in U':f(w,z)=0\}=\{(w,g(w))\in U:w\in W\}
	\end{equation*}
\end{theorem}



\begin{propo}\label{biholomorfismojac}
	Sea $f:U\to V$ una aplicación holomorfa biyectiva entre dos subconjuntos abiertos $U,V\subset \cc^n$. Entonces para todo $z\in U$ se tiene que $\det J(f)(z)\neq 0$. En particular, $f$ es un biholomorfismo.
\end{propo}

\section{Comportamiento local}\label{complocal}
	
	En una variable compleja, el comportamiento local de las funciones holomorfas es muy simple. Si $f:U\to\cc$ es una función holomorfa en alguna vecindad $U\subset\cc$ del cero, entonces $f$ se escribe de forma única como
	\begin{equation*}
		f(z)=u(z)z^k
	\end{equation*} 
	para algún $k\in\nn$ y una función $u\in\mathscr{O}_n(U)$ unidad. Esto significa que $u(0)\neq 0$ y $1/u(z)$ está bien definida en alguna vecindad del cero. En varias variables complejas la situación es mucho más complicada.\\
	\begin{definition}
		Sean $U,V\subset \cc^n$ vecindades del $0\in\cc^n$. Decimos que $f\in\mathscr{O}_n(U)$ y $g\in\mathscr{O}_n(V)$ están relacionadas si existe una vecindad $W\subset\cc^n$ del $0$ tal que
		\begin{equation*}
			f|_{U\cap W}=g|_{V\cap W}
		\end{equation*} 
		Llamamos a la clase de equivalencia $[f]$ \textbf{germen} de la función $f$ en el cero.
		Denotamos por $\oo_{n}$ al conjunto de gérmenes de funciones en el cero.
	\end{definition}
	
	\begin{remark}
		Si $f,g\in \oo_{n}$ son gérmenes, entonces podemos definir las operaciones
		\begin{align*}
			[f]+[g]:=[f+g]\\
			[f][g]:=[fg].
		\end{align*}
		Por lo tanto, el conjunto $\oo_{n}$ es un anillo. Más aún, $\cc$ está contenido en $\oo_{n}$ a travez de las funciones constantes y este anillo puede ser descrito por el límite directo
		\begin{equation*}
			\oo_{n}=\lim_{0\in U}\oo_n(U).
		\end{equation*} 
	\end{remark}
	
	Nosotros podemos pensar a un germen $[f]\in\oo_{n}$ como una función holomorfa alrededor del cero en alguna vecidad no especificada. Más aún, el valor $f(0)\in\cc$ está bien definido para gérmenes, pero no está bien definidido para cualquier otro punto.\\
	
	\begin{remark}
		Por definición, una función $f$ es holomorfa en $0\in\cc^n$ si puede ser expandido como una serie de potencias convergente 
		\begin{equation*}
			f(z)=\sum_{\alpha\in\nn}a_\alpha z^\alpha.
		\end{equation*}
		De está manera, tenemos un isomorfismo de anillos 
		\begin{equation*}
			\oo_{n}\cong \cc\{z_1,\cdots,z_n\}
		\end{equation*}
		donde $\cc\{z_1,\cdots,z_n\}$ es el anillo de series convergentes (alrededor del cero) en las variables $z_1,\cdots,z_n$.
	\end{remark}

	\begin{example}
		Para $n=0$, tenemos que $\oo_{n}\cong\cc$. Para $n=1$, se cumple que $\oo_{n}\cong\cc\{z\}$. Así, las funciones holomorfas alrededor del cero en una variable corresponden algebraicamente con el anillo $\cc\{z\}$. Este anillo es un anillo de valuación discreta, esto es, todo ideal de $\cc\{z\}$ es de la forma $\langle z^k \rangle$ para algún $k\in\nn$.
	\end{example}
	
	De esta manera, estudiar el comportamiento local de las funciones holomorfas es equivalente a estudiar las propiedades algebraicas del anillo $\oo_{n}$.
	
	\begin{remark}[$\oo_{n}$ es un anillo local]
		Consideremos el conjunto de elementos no invertibles en $\oo_{n}$
		\begin{equation*}
			\mathfrak{m}_n=\{f\in\oo_{n}:f(0)=0\}.
		\end{equation*}
		Ya que si $f(0)\neq 0$, entonces $1/f$ es holomorfa en una vecindad alrededor del origen. Así, $1/f\in\oo_{n}$. Más aún, $\mathfrak{m}_n$ es un ideal pues si $f\in\oo_{n}$ y $m\in\mathfrak{m}_n$, entonces $(fm)(0)=0$. Con esto $fm\in\oo_{n}$, es decir, $\mathfrak{m}_n$ es un ideal de $\oo_{n}$. Esto significa que $\oo_{n}$ es un anillo local.\\
		En particular, si $f\in\mathfrak{m}_n\subset \cc\{z_1,\cdots,z_n\}$ entonces su expansión en serie no tiene coeficiente independiente
		\begin{equation*}
			f=\sum_{\alpha\in\nn^n\setminus\{0\}}a_\alpha z^{\alpha}.
		\end{equation*}
		Esto significa que el ideal $\mathfrak{m}_n$ es finitamente generado y
		\begin{equation*}
			\mathfrak{m}_n=\langle z_1,\dots,z_n \rangle.
		\end{equation*}
	\end{remark}
	Notemos que el campo residual $\oo_{N}/\mathfrak{m}_n$ es isomorfo a $\cc$. En particular, el entero $n$ (la dimensión de $\cc^n$) se puede recuperar del anillo local $\oo_{n}$. Dado que el ideal maximal $\mathfrak{m}_n$ está generado por lo elementos $z_1,\dots,z_n$, entonces imagen de los generadores bajo el cociente
	\[\begin{tikzcd}
		{\mathfrak{m}_n} && {\mathfrak{m}_n/\mathfrak{m}_n^2}
		\arrow["\rho", from=1-1, to=1-3]
	\end{tikzcd}\]
	genera una base para $\mathfrak{m}_n/\mathfrak{m}_n^2$. Por lo que:
	\begin{corollary}
		El cociente $\mathfrak{m}_n/\mathfrak{m}_n^2$ es un $\cc$-espacio vectorial de dimensión $n$.
	\end{corollary}
	
	Otra propiedad importante y simple de demostrar de este anillo es el siguiente resultado.
	
	\begin{propo}
		El anillo $\oo_{n}$ es un dominio entero.
	\end{propo}
	
	Este anillo tiene propiedades aún más profundas e importantes, pero tenemos que estudiar la estructura local de las funciones holomorfas con más cuidado. Notemos que tenemos una inclusión natural de anillos 
	\begin{equation*}
		\oo_{n-1}\subset \oo_{n-1}[z_n]\subset \oo_n.
	\end{equation*}
	Los elementos del anillo $\oo_{n-1}[z_n]$ son polinomios de la forma $z_n^k+a_1z_n^{k-1}+\cdots+a_k$ con $a_1,\dots,a_k\in \oo_{n-1}$, los cuales son gérmenes en $\oo_n$. Además, la primera inclusión tiene una naturaleza algebraica mientras que la segunda es más analítica. Ahora vamos a presentar el resultado fundamental para entender esta segunda inclusión llamado \textit{Teorema de Preparación de Weierstraß}.
	
	\begin{definition}
		Un elemento $z_n^k+a_1z_n^{k-1}+\cdots+a_k\in\oo_{n-1}[z_n]$ con $k\geq 1$ es llamado \textbf{polinomio de Weierstraß} si $a_1,\dots,a_k\in \mathfrak{m}_{n-1}$.
	\end{definition}
		Análogamente al caso de una varieable, queremos ver como esencialmente todo germen $f\in\oo_n$ puede ser escrito de la forma $f=uh$ donde $h$ es un polinomio de Weierstraß y $u$ es unidad. Sin embargo, como $h(0,z_n)=z_n^d$, tenemos que si $f=uh$ entonces $f$ restringida a la línea $w=(z_1,\dots,z_{n-1})$ no puede ser identicamente cero.  
	\begin{definition}
		Sea $U\subset\cc^n$ una vecindad del cero y $f\in\oo_n(U)$. Decimos $f$ es \textbf{regular} en $z_n$ si la función holomorfa $f(0,z_n)$ no es identicamente cero. Además, como $f(0,z_n)=u(z_n)z_n^d$ para algún $d$, decimos que $f$ es regular de \textbf{orden} $d$.
	\end{definition}
		La notación de regularidad tiene sentido para los elementos de $\oo_n$ ya que solo depende del comportamiento de $f$ en alguna vecindad albitrariamente pequeña del origen.
		
	\begin{theorem}[Teorema de Preparación de Weierstraß]
		Si $f\in\oo_n$ es regular de orden $d$ en $z_n$, entonces existe un único polinomio de Weierstraß $h\in\oo_{n-1}[z_n]$ de grado $d$ tal que $f=uh$ para alguna unidad $u\in\oo_n$.
	\end{theorem}
		La idea de la demostración es un poco simple: Fijando $w\in\cc^{n-1}$ suficientemente cerca del origen, y consideremos $f_w(z_n):=f(w,z_n)$ la cual es una función holomorfa en $z_n$. Ya que $f_0(z_n)$ tiene un cero de orden $d$, entonces cada $f_w(z_n)$ tiene exactamente $d$ ceros (contando sus multiplicidades) cerca del origen. Escribimos como $a_1(w),\dots,a_d(w)$ las raíces y se define la función $u_w(z_n)=(z_n-a_1(w))\cdots(z_n-a_d(w))$. El objetivo principal de la prueba es demostrar que la función $h=f_w(z_n)/u_w(z_n)$ es holomorfa en $w$.\\
		La prueba formal y completa se puede consultar en \parencite[Cap. 1, Sec, 1.1, Prop. 1.1.6]{huybrechts2005complex}.
		
	\begin{propo}\label{anillodfu}
		El anillo $\oo_{n}$ es un dominio de factorización única.
	\end{propo}
	
	Existe una manera de hacer división larga de polinomios de Weierstraß, al igual que el anillo $\cc[z]$ y este resultado también es fundamental sobre el álgebra local.
	
	\begin{theorem}[Teorema de división de Weierstraß]
		Sea $h\in\oo_{n-1}[z_n]$ un polinomio de Weierstraß de grado $d$. Entonces para todo $f\in\oo_n$ puede ser escrito de manera única de la forma
		\begin{equation*}
			f=qh+r,\qquad\text{para }q\in \oo_n 
		\end{equation*} 
		y $r\in\oo_{n-1}[z_n]$ es un polinomio de grado $<d$. Más aún. si $f\in\oo_{n-1}[z_n]$, entonces también $q\in\oo_{n-1}[z_n]$.
	\end{theorem}
	\begin{proof}
		\parencite[Cap. 1, Sec, 1.1, Prop. 1.1.17]{huybrechts2005complex}
	\end{proof}
	
	De esta manera tenemos nuestra última propiedad importante del anillo de gérmenes de funciones holomorfas.
	
	\begin{theorem}\label{anillonoetheriano}
		El anillo $\oo_{n}$ es un anillo Noetheriano.
	\end{theorem}
	\begin{proof}
		Vamos a probar el teorema por indución sobre $n$. El caso para $n=0$ es trivial, pues el anillo $\oo_n=\cc$ es noetheriano.\\
		Supongamos que $\oo_{n-1}$ es un anillo noetheriano. Entonces por el teorema de la base de Hilbert, el anillo de polinomios $\oo_{n-1}[z_1]$ es noetheriano. Sea $I\subset\oo_n$ un ideal no cero y tomemos $0\neq f\in I$. Haciendo un cambio de coordenadas de ser necesario, podemos aplicar el teorema de Preparación de Weierstraß a $f$, es decir, existe un polinomio de Weierstraß $g\in \oo_{n-1}[z_1]$ y $h\in\oo_n$ unidad tal que
		\begin{equation*}
			f=gh.
		\end{equation*}
		En particular, $g\in I$. Luego, el ideal $I\cap \oo_{n-1}[z_1]$ es finitamente generado en $\oo_{n-1}[z_1]$, digamos
		\begin{equation*}
			I\cap \oo_{n-1}[z_1]=\langle g_1,\dots,g_k \rangle.
		\end{equation*}
		Para cualquier otro $\tilde{f}\in I$, el teorema de división de Weierstraß nos dice que 
		\begin{equation*}
			\tilde{f}=g\tilde{h}+r
		\end{equation*}
		para algún $r\in\oo_{n-1}[z_1]$. Más aún, como $\tilde{f},g,\tilde{h}\in I$, también $r\in I$ y podemos concluir que $r\in I\cap \oo_{n-1}[z_1]$. De esta manera,
		\begin{equation*}
			\tilde{f}=g\tilde{h}+\sum_{i=1}^{k}a_ig_i.
		\end{equation*}
		Por lo tanto, $I$ está generado por los elementos $\{g,g_1,\dots,g_k\}$.
	\end{proof}
	
	
\section{Álgebra de los conjuntos analíticos}
	
		Usaremos todas las propiedades algebraicas del anillo local $\oo_n$ para describir el álgebra y la geometría los conjuntos analíticos en $\cc^n$. En primer lugar, definamos una relación de equivalencia sobre subconjuntos de $\cc^n$.
	\begin{definition}
			Sean $X,Y\subset \cc^n$ y $a\in X\cap Y$.
			\begin{enumerate}
				\item Decimos  que $X$ y $Y$ son \textbf{equivalentes en} $a$ si existe una vecindad abierta $U$ de $a$ tal que
				\begin{equation*}
					X\cap U = Y\cap U.
				\end{equation*}
				\item Una clase de equivalencia de $X$ es  llamada \textbf{germen} de $X$ en $a$, y es denotada por
				\begin{equation*}
					[X_a].
				\end{equation*}
				Se omite el punto $a$ cuando el punto es fijo o claro.
				\item Si $[X]$ y $[Y]$ son gérmenes de conjuntos, definimos
				\begin{align*}
					[X]\cap[Y]&:=[X\cap Y]\\
					[X]\cup[Y]&:=[X\cup Y]\\
					[X]\subset[Y]&:\Longleftrightarrow [X]\cap[Y]=[X].
				\end{align*}
				\item Si $[f]\in\oo_0$ es representado por algún $f\in\oo(U)$ donde $U\subset\cc^n$ es alguna vecindad de origen, definimos
				\begin{equation}
					Z([f]):=[Z(f)].
				\end{equation}
				\item [5. ]Si $[X_0]$ es un germen del cero y $[f]\in\oo_0$, decimos que $[f]$ se anula en $[X]$ si
				\begin{equation*}
					[X]\subset Z([f]).
				\end{equation*}
			\end{enumerate}
	\end{definition}

	Nos centraremos en estudiar conjuntos analíticos solo en el origen, por lo que el concepto de germen es necesario para solo conservar su información alrededor del origen.
	
	\begin{definition}
			Un germen $[X]$ es llamado \textbf{germen analítico} o \textbf{conjunto analítico} si existe una cantidad finita de gérmenes de funciones holomorfas $[f_1],\cdots,[f_k]\in\oo_0$ tal que
			\begin{equation*}
				[X]=Z([f_1],\cdots,[f_k])
			\end{equation*}
	\end{definition}
		
	\begin{propo}
			Si $[X]$ y $[Y]$  son gérmenes analíticos,  entonces $[X]\cap [Y]$ y $[X]\cup [Y]$ también son gérmenes analíticos.
	\end{propo}
	
		\begin{definition}
			Sea $[X]$ un germen analítico en el cero y $A\subset\oo_n$. Definimos
			\begin{equation*}
				I([X]):=\{[f]\in\oo_0:[X]\subset Z([f])\}
			\end{equation*}
			y
			\begin{equation*}
				Z(A)=\bigcap_{[f]\in A}Z([f]).
			\end{equation*}
		\end{definition}
		
	\begin{example}
		Calculando el $I$ del germen del conjunto $\{0\}$,
		\begin{align*}
			I([\{0\}])&=\{[f]\in\oo_0:[f]\text{ se anula  en }[\{0\}]\}\\
			&=\{f\text{ holomorfa en el cero}:f(0)=0\}\\
			&=\mathfrak{m}_n,
		\end{align*}
		el cual es el ideal maximal de $\oo_0$.
		Ahora, calculando los ceros de todo el anillo $\oo_0$,
		\begin{equation*}
			Z(\oo_n)=\bigcap_{[f]\in\oo_0}Z([f])=\emptyset
		\end{equation*}
		ya que $\oo_0$ contiene las constantes.
	\end{example}
	
	
	\begin{propo} $ $
		\begin{enumerate}
			\item El conjunto $I([X])$ es un ideal del anillo $\oo_n$, llamado \textbf{ideal asociado}.
			\item El conjunto $Z(A)=Z(\langle A\rangle)$ es  un germen analítico.
		\end{enumerate}
	\end{propo}
	
	\begin{propo}
		Sea $[X]$ un germen analítico. Entonces
		\begin{equation*}
			Z(I([X]))=[X].
		\end{equation*}
	\end{propo}
	
	\begin{theorem}[de Rückert]
		Sea $\mathfrak{a}\subset\oo_0$ cualquier ideal. Entonces
		\begin{equation*}
				I(Z([\mathfrak{a}]))=\sqrt{\mathfrak{a}}.
		\end{equation*}
	\end{theorem}
	
	\begin{remark}
		Consideremos un ideal radical $J\subset\oo_0$, entonces existe un ideal $\mathfrak{a}\subset\oo_0$ tal que
		\begin{equation*}
			J=\sqrt{\mathfrak{a}}
		\end{equation*}
		Luego, por el teorema de Rückert
		\begin{align*}
			I(Z(J))=\sqrt{J}=\sqrt{\sqrt{\mathfrak{a}}}=\sqrt{\mathfrak{a}}=J
		\end{align*}
		Esto nos da una relación biyectiva entre los dos conjuntos
		\begin{align*}
			\begin{Bmatrix}
				\text{Gérmenes analíticos}\\
				\text{de }\cc^n
			\end{Bmatrix}\longleftrightarrow\begin{Bmatrix}
				\text{Ideales radicales}\\
				\text{de }\cc\{z_1,\dots,z_n\}
			\end{Bmatrix}
		\end{align*}
		Más aún, tenemos la relación
		\begin{align*}
			\begin{Bmatrix}
				\text{Gérmenes analíticos}\\
				\text{irreducibles de }\cc^n
			\end{Bmatrix}\longleftrightarrow\begin{Bmatrix}
				\text{Ideales primos}\\
				\text{de }\cc\{z_1,\dots,z_n\}
			\end{Bmatrix}
		\end{align*}
	\end{remark}
	
	\begin{definition}
		Sea $X$ un conjunto analítico irreducible definido por un ideal primo $\mathfrak{p}\subset\oo_n$. Definimos la \textbf{dimensión} de $X$ como
		\begin{equation*}
			\dim X:=\mathrm{dim}_{\text{krull}}\oo_n/\mathfrak{p}.
		\end{equation*}
		Además, definimos la \textbf{codimensión} de $X$ como $\mathrm{codim}X:=n-\dim X$.
	\end{definition}
	
	Sea $X\subset\cc^n$ un germen analítico irreducible de codimensión $1$, es decir, $\dim X=n-1$. Entonces el ideal primo $\mathfrak{p}$ asociado a $X$ tiene altura $1$ en $\oo_n$. Un resultado esencial de los DFU es que todo ideal de altura $1$ es principal. Esto significa que 
	\begin{equation*}
		\mathfrak{p}=\langle f \rangle 
	\end{equation*}
	para algún elemento irreducible $f\in\oo_n$. Por lo tanto, un germen analítico irreducible de codimensión uno está definido por una función holomorfa irreducible.
	
	\begin{propo}\label{hipersuperficiesenCn}
		Todo germen analítico irreducible de codimension uno es el conjunto de ceros de una sola función holomorfa irreducible.
	\end{propo}

	